\section{Scope Boundaries}
\label{sec:scope}

This section defines explicit boundaries for Timeline Compiler. Clarity here prevents scope creep and ensures the system remains focused on its core value proposition: transforming structured plan descriptions into presentation-quality visual timelines.

\subsection{Governing Principle}

The scope boundary is defined by the \emph{rendering abstraction}: Timeline Compiler transforms a communication-oriented intermediate representation into visual output. It does not generate, compute, or manage the plan data itself.

\begin{quote}
\emph{If a feature requires understanding the semantics of work — dependencies, resources, costs, velocity — it is out of scope. If a feature improves the visual communication of a plan already decided, it is potentially in scope.}
\end{quote}

\subsection{In Scope}

The following capabilities are within Timeline Compiler's scope:

\subsubsection{Visual Artifact Types}

\begin{description}
  \item[Timelines] Horizontal time-based visualizations showing activities, milestones, and phases across a date range.
  
  \item[Roadmaps] Product or program roadmaps with tracks (swimlanes), grouped activities, and strategic milestones.
  
  \item[Milestone Views] Focused views highlighting key dates and deliverables without detailed activity bars.
  
  \item[Phase Summaries] High-level views showing phases or stages of a program without granular activities.
  
  \item[Swimlane Diagrams] Multi-track views where parallel workstreams are visually separated.
\end{description}

\subsubsection{IR Elements}

The IR supports the following semantic elements (detailed in Section~\ref{sec:ir}):

\begin{description}
  \item[\texttt{tracks}] Named swimlanes or rows that group related activities. Tracks provide visual and semantic organization.
  
  \item[\texttt{groups}] Logical groupings within or across tracks for nested organization (e.g., workstreams within a track).
  
  \item[\texttt{activities}] Time-bounded items (bars) representing work, initiatives, or phases. Activities have start/end dates, labels, and optional metadata.
  
  \item[\texttt{milestones}] Point-in-time markers representing key dates, deliverables, or decisions.
  
  \item[\texttt{sections}] Temporal divisions of the timeline (e.g., quarters, phases) that span all tracks.
  
  \item[\texttt{range}] The overall time window for the visualization.
  
  \item[\texttt{date ranges}] Start and end dates for activities and sections. Supports dates, months, quarters, and years.
  
  \item[\texttt{progress}] Numeric (0.0–1.0) indication of completion for activities. Purely visual — not computed.
  
  \item[\texttt{status}] Semantic status indicators (e.g., \texttt{planned}, \texttt{in-progress}, \texttt{complete}, \texttt{at-risk}, \texttt{blocked}).
  
  \item[\texttt{labels}] Text labels for all elements. Brief, presentation-appropriate text.
  
  \item[\texttt{annotations}] Callouts, notes, or emphasis markers attached to dates or elements.
  
  \item[\texttt{legends}] Explanatory legends for status colors, track semantics, or custom indicators.
\end{description}

\subsubsection{Presentation Features}

\begin{description}
  \item[Theming] Complete visual customization via theme files: colors, typography, spacing, shapes.
  
  \item[Emphasis] Visual mechanisms to highlight specific activities or milestones (e.g., bold, glow, callout).
  
  \item[Today Markers] Visual indicator for the current date (provided as input, not computed).
  
  \item[Progress Visualization] Visual representation of progress within activity bars.
  
  \item[Status Indicators] Color or icon mapping for semantic statuses.
  
  \item[Annotations] Callouts, arrows, or labels that highlight specific points.
\end{description}

\subsubsection{Output Formats}

\begin{description}
  \item[SVG] Primary vector output for web embedding and further processing.
  
  \item[PNG] Rasterized output for contexts requiring bitmap images.
  
  \item[PDF] Print-quality vector output for documents and decks.
  
  \item[PPTX] PowerPoint-compatible output for direct use in presentations.
\end{description}

\subsubsection{Tooling}

\begin{description}
  \item[CLI] Command-line interface for rendering (\texttt{timeline render input.yaml -o output.pdf}).
  
  \item[Library API] Programmatic access for embedding in other tools.
  
  \item[Schema Validation] Validation of IR against formal schema.
  
  \item[Theme Validation] Validation of theme files.
\end{description}

\subsection{Out of Scope}

The following capabilities are explicitly excluded:

\subsubsection{Project Management Semantics}

\begin{description}
  \item[Dependency Scheduling] Timeline Compiler does not compute dates from dependencies. If Activity B depends on Activity A, the author must provide both dates explicitly. The IR may optionally \emph{record} dependencies for visual rendering (e.g., arrows), but it does not \emph{resolve} them.
  
  \textbf{Rationale:} Dependency resolution is core project-management logic. It requires understanding task semantics, durations, constraints, and often iterative solving. This is the domain of MS Project, Primavera, and similar tools. Timeline Compiler consumes the \emph{output} of scheduling, not the inputs.
  
  \item[Resource Management] No support for assigning people, teams, or resource pools to activities. No resource leveling or capacity planning.
  
  \textbf{Rationale:} Resource management requires organizational context (who is available, at what capacity, with what skills) that is outside the scope of a rendering compiler.
  
  \item[Critical Path Analysis] No calculation of critical paths or slack time.
  
  \textbf{Rationale:} Critical path analysis is a scheduling concern, not a rendering concern. Upstream tools compute this; Timeline Compiler may visualize the result if provided.
  
  \item[Sprint/Iteration Tracking] No support for agile sprints, velocity tracking, burndowns, or iteration planning.
  
  \textbf{Rationale:} Sprint tracking is operational work management. Timeline Compiler visualizes strategic communication, not operational tracking.
  
  \item[Cost Management] No budgets, cost tracking, earned value, or financial projections.
  
  \textbf{Rationale:} Financial data requires integration with accounting and project-control systems. This is far outside the rendering scope.
  
  \item[Portfolio Optimization] No multi-project prioritization, resource allocation across projects, or strategic portfolio balancing.
  
  \textbf{Rationale:} Portfolio management is a strategic planning function, not a visualization function.
  
  \item[Baseline Comparison] No tracking of baseline vs.\ actual schedules over time.
  
  \textbf{Rationale:} Baseline tracking requires temporal versioning and comparison logic. This is project-control functionality.
  
  \item[Risk Registers] No formal risk tracking, probability/impact matrices, or mitigation planning.
  
  \textbf{Rationale:} Risk management is a project-management discipline with its own tooling requirements.
\end{description}

\subsubsection{Data Management}

\begin{description}
  \item[Persistent Storage] Timeline Compiler does not store timelines. It transforms input files to output files. Storage is the user's responsibility.
  
  \item[Collaboration Features] No real-time collaboration, comments, change tracking, or multi-user editing.
  
  \item[Version History] No built-in versioning. Users should use Git or similar tools.
\end{description}

\subsubsection{Interactive Features}

\begin{description}
  \item[Clickable Outputs] Output formats (SVG, PNG, PDF, PPTX) are static. Hyperlinking or interactive drill-down is out of scope for MVP, though SVG hyperlinks may be a future extension.
  
  \item[WYSIWYG Editing] No drag-and-drop authoring. Timeline Compiler is a compiler, not an editor.
  
  \item[Live Preview During Authoring] Out of scope for core; a VS Code extension may provide this separately.
\end{description}

\subsubsection{Data Ingestion}

\begin{description}
  \item[Native Integrations] Timeline Compiler does not natively read from Jira, Azure DevOps, Asana, GitHub Projects, or other systems. It consumes only its IR format.
  
  \textbf{Note:} Separate adapter tools or scripts may transform external data to the IR, but these are outside Timeline Compiler's core scope.
\end{description}

\subsection{Boundary Cases}

Some features lie at the boundary and require explicit decisions:

\subsubsection{Dependency Arrows (Visual Only)}

\textbf{Decision:} IN SCOPE — visual only.

The IR may include dependency declarations for the purpose of rendering arrows or lines between activities. However, dependencies are purely visual annotations; they do not affect date computation or validation.

\begin{lstlisting}[language=yaml,caption={Dependencies as visual annotations}]
activities:
  - id: design
    label: "Design Phase"
    range: { start: 2026-01, end: 2026-03 }
  - id: build
    label: "Build Phase"
    range: { start: 2026-04, end: 2026-08 }
    depends_on: [design]  # Visual arrow only
\end{lstlisting}

\subsubsection{Today Line}

\textbf{Decision:} IN SCOPE — provided as input.

The ``today'' date may be provided in the IR (e.g., \texttt{today: 2026-06-09}) and rendered as a vertical marker. The compiler does not infer today's date from the system clock; determinism requires all inputs to be explicit.

\subsubsection{Auto-Layout}

\textbf{Decision:} IN SCOPE — deterministic algorithms only.

The compiler automatically lays out activities within tracks (e.g., stacking overlapping items). Layout algorithms must be deterministic. No randomized or heuristic placement that could vary between runs.

\subsubsection{External Image Embedding}

\textbf{Decision:} DEFERRED to future.

Embedding external images (logos, icons) in output is not in MVP scope. Theme-provided icons (e.g., milestone diamonds) are in scope.

\subsection{Scope Summary}

\begin{table}[h]
\centering
\caption{Scope Boundaries Summary}
\begin{tabularx}{\textwidth}{lX}
\toprule
\textbf{Category} & \textbf{Scope Status} \\
\midrule
Timelines, roadmaps, milestone views & In Scope \\
Tracks, activities, milestones, sections, annotations & In Scope \\
Progress and status visualization & In Scope \\
Theming and visual customization & In Scope \\
SVG, PNG, PDF, PPTX output & In Scope \\
CLI and library API & In Scope \\
Dependency arrows (visual only) & In Scope \\
Today marker (explicit input) & In Scope \\
\midrule
Dependency scheduling/resolution & Out of Scope \\
Resource management & Out of Scope \\
Critical path analysis & Out of Scope \\
Sprint/iteration tracking & Out of Scope \\
Cost management & Out of Scope \\
Portfolio optimization & Out of Scope \\
Baseline comparison & Out of Scope \\
Risk registers & Out of Scope \\
Data storage and collaboration & Out of Scope \\
Interactive/clickable output & Out of Scope \\
WYSIWYG editing & Out of Scope \\
Native data source integrations & Out of Scope \\
\bottomrule
\end{tabularx}
\end{table}
